\documentclass[12pt,letterpaper]{article}
\usepackage[utf8]{inputenc}
\usepackage[T1]{fontenc}
\usepackage{lmodern,amsmath,amssymb,array,enumitem,xcolor,geometry,tcolorbox}
\geometry{left=1.7cm,right=1.7cm,top=1.5cm,bottom=1.6cm}
\setlength{\parindent}{0pt}
\definecolor{azul}{HTML}{17365D}
\definecolor{claro}{HTML}{EAF1F8}
\newcommand{\familia}[2]{\section*{Familia #1. #2}}
\newcommand{\campo}[1]{\textcolor{azul}{\textbf{#1}}}
\begin{document}
\begin{center}
{\LARGE\bfseries\color{azul} SUCESIONES Y LÍMITES}\par
{\Large Banco paramétrico -- Nivel medio}
\end{center}
\begin{tcolorbox}[colback=claro,colframe=azul]
Nivel de referencia de la primera guía. Cada pregunta exige cálculo y una justificación
breve. Las tuplas fueron elegidas para evitar operaciones extensas y mantener resultados
exactos. No intercambiar componentes entre tuplas dependientes.
\end{tcolorbox}

\familia{1}{Patrón alternante y cálculo de términos}
\campo{Plantilla:}
\[
 a_n=(-1)^{n+h}\frac{pn+q}{n+r},\qquad h\in\{0,1\}.
\]
Pedir cuatro términos, uno posterior y explicar el efecto de $h$.\par
\campo{Condiciones:} $p\in\{2,3,4\}$, $|q|\leq4$, $r\in\{1,2,3,4\}$,
$pn+q>0$ para $n\geq1$ y $5\leq k\leq10$.\par
\campo{Tuplas $(p,q,r,h,k)$:}
\[
 (2,3,4,1,8),\ (3,1,2,0,7),\ (4,-1,3,1,6),\ (2,1,1,0,9).
\]

\familia{2}{Conversión explícita-recursiva}
\campo{Plantilla:}
\[
 b_n=L+Cq^{n-1}.
\]
Pedir una recurrencia de primer orden y comprobarla.\par
\campo{Condiciones:} $q\in\{2,3,-2\}$, $C\neq0$, $|C|\leq4$, $|L|\leq5$;
usar
\[
 b_{n+1}=q b_n+(1-q)L,\qquad b_1=L+C.
\]
\campo{Tuplas $(L,C,q)$:}
\[
 (3,2,2),\ (-1,3,2),\ (2,-1,3),\ (1,2,-2).
\]

\familia{3}{Monotonía racional}
\campo{Plantilla:}
\[
 c_n=\frac{An+B}{Cn+D}.
\]
Pedir determinar la monotonía mediante diferencia.\par
\campo{Condiciones:} $C>0$, $Cn+D>0$ para $n\geq1$, $AD-BC\neq0$ y
$|A|,|B|,|C|,|D|\leq6$.\par
\campo{Control:}
\[
 c_{n+1}-c_n=\frac{AD-BC}{(Cn+D)(Cn+C+D)}.
\]
\campo{Tuplas $(A,B,C,D)$:}
\[
 (4,-1,2,3),\ (3,2,2,5),\ (2,5,3,1),\ (5,1,2,4).
\]

\familia{4}{Acotación por una asíntota horizontal}
\campo{Plantilla:}
\[
 d_n=L+\frac{K}{n+r}.
\]
Pedir monotonía, una cota alcanzada y otra no alcanzada.\par
\campo{Condiciones:} $K\neq0$, $r\geq0$ y $n+r>0$. Si $K>0$, decrece hacia
$L$; si $K<0$, crece hacia $L$.\par
\campo{Tuplas $(L,K,r)$:}
\[
 (3,4,1),\ (5,-3,2),\ (-1,2,1),\ (2,-4,3).
\]
\campo{Control:} extremo alcanzado $d_1=L+K/(r+1)$; límite $L$.

\familia{5}{Límite racionalizable}
\campo{Plantilla:}
\[
 e_n=\sqrt{p^2n^2+qn+r}-pn.
\]
Pedir racionalizar y calcular el límite.\par
\campo{Condiciones:} $p\in\{1,2,3,4\}$, $q$ múltiplo de $2p$ para obtener
$q/(2p)$ entero o semientero, y el radicando debe ser positivo para $n\geq1$.\par
\campo{Tuplas $(p,q,r)$:}
\[
 (1,4,1),\ (2,4,3),\ (3,6,2),\ (4,8,1).
\]
\campo{Control:} $\displaystyle\lim e_n=q/(2p)$.

\familia{6}{Recurrencia afín convergente}
\campo{Plantilla:}
\[
 x_1=u,\qquad x_{n+1}=\rho x_n+(1-\rho)L.
\]
Pedir demostrar monotonía, acotación y convergencia.\par
\campo{Condiciones:} $0<\rho<1$; usar $u<L$ para crecimiento y $u>L$ para
decrecimiento.\par
\campo{Tuplas $(\rho,L,u)$:}
\[
 \left(\tfrac12,4,0\right),\
 \left(\tfrac13,3,0\right),\
 \left(\tfrac14,2,6\right),\
 \left(\tfrac23,5,8\right).
\]
\campo{Control:} $x_n=L+(u-L)\rho^{n-1}$ y $\lim x_n=L$.

\familia{7}{Teorema del sándwich}
\campo{Plantilla:}
\[
 y_n=\varepsilon_n\frac{an+b}{n^2+c},
\qquad |\varepsilon_n|\leq1.
\]
Usar $\varepsilon_n=(-1)^n$, $\sin n$ o $\cos(n^2)$ y pedir justificar el límite.\par
\campo{Condiciones:} $a>0$, $b,c\geq0$, $a,b,c\leq6$.\par
\campo{Tuplas $(\varepsilon_n,a,b,c)$:}
\[
 ((-1)^n,2,3,1),\ (\sin n,3,1,2),\ (\cos(n^2),4,2,3),\ ((-1)^{n+1},5,1,4).
\]
\campo{Control:}
\[
 |y_n|\leq\frac{an+b}{n^2+c}\longrightarrow0.
\]

\end{document}
